\begin{center}
{\Large\bfseries A vanishing theorem on the rank-three host}\\[5pt]
{\large With the finite data for spectral family selection}\\[4pt]
Edition 0.60 \quad 24 November 2026
\end{center}

The host used in the rank-three counting theorem comes with an
arithmetic character and prescribed local conditions. Its
weight $-2$ Selmer group measures whether a locally admissible
cocycle is globally a coboundary. We prove that this group is zero.
The proof begins with the Dirichlet characters to which the Bernoulli
vanishing theorem applies. It then identifies the host module with a
line, checks the local conditions on its cocycles at all places,
and returns the global coboundary to the given host frame.
In the project's dependency register this closes the shared
RH and Dirichlet-GRH obligation \texttt{TASK\_HOST\_XPP\_INSTALL}.

The spectral chapters continue the period identity of edition 0.59.
They treat the finite fibres, dependence on the choice of a term
representative, Haar measures and the analytic data needed for one
family of spectral settings. The local and family hypotheses for each
result are stated where it is used.

\begingroup\emergencystretch=2em
\section{Constructing the Bernoulli input on the host}
\subsection{The two Dirichlet characters attached to a host}
\label{sec:host-xpp-character-input}

% BEGIN HXD1
\paragraph{HXD1. Splitting the actual host character.}
Fix the host's algebraic closure $\Omega$ of $\mathbb Q$, its odd
prime $\ell$, its coefficient DVR $O=c.O$, its level $N>0$, its
exponent $e$, and its Dirichlet character $\psi$ of level $N$.
Write $k=O/\mathfrak m_O$. The ring $O$ is finite over
$\mathbb Z_\ell$ and contained in $\overline{\mathbb Q}_\ell$.
Its residue field is finite of characteristic $\ell$. Indeed, a
maximal ideal of a finite integral extension contracts to a maximal
ideal, so $\ell O\subseteq\mathfrak m_O$, and $O/\ell O$ is finite
over $\mathbb F_\ell$. In this section a character evaluated on a
unit is identified with its unit value, and the bar denotes reduction
to $k$.

Split $\psi=\theta_0\kappa$ into its prime-to-$\ell$ and
$\ell$-power parts in the finite cyclic group generated by $\psi$.
In the CRT decomposition of $(\mathbb Z/N)^\times$ let $\psi_2$ be
the component on the factor at $2$, inflated to level $N$. That factor
is a $2$-group, hence $\kappa$ is trivial on it. The $2$-component
of $\theta_0$ equals $\psi_2$. Thus the source character
$\theta=\operatorname{th1}(\ell,\psi)$, obtained by dividing
$\theta_0$ by its $2$-component, satisfies
\[
     \psi=\psi_2\theta\kappa,\qquad
     \ell\nmid\operatorname{ord}(\theta),\qquad
     \operatorname{cond}(\theta^2)\text{ is odd}.
\]
The conductor assertion follows because $\theta$ factors through the
odd part of $N$. Reduction sends $\kappa$ to $1$ on all units, since
$X^{\ell^s}-1=(X-1)^{\ell^s}$ in $k[X]$ for any $\ell^s$ killing
$\kappa$. Put
\[
 t=c_2(\psi),\quad n=t-1,\quad
 d=\operatorname{ord}_{\operatorname{lcm}(4,\operatorname{ord}\theta)}(\ell),
 \qquad L=\ell^d-1.
\]
Here $\operatorname{cond}(\psi_2)=2^t$ and $2^t\mid N$.
The multiplicative order exists since $\ell$ is coprime to the
modulus. Therefore $d>0$, $4\mid L$, $\operatorname{ord}(\theta)\mid L$,
and $\ell-1\mid L$ by the geometric sum formula.
The host's $n_0(\ell,\psi)$ equals $v_2(L)$ by definition.
% END HXD1

% BEGIN HXD2
\paragraph{HXD2. The primitive pair and the conductor bound.}
Use the Teichm\"uller character $\omega:(\mathbb Z/\ell)^\times
\to\mathbb Z_\ell^\times\to O^\times$. The unique lift of each
nonzero residue to a root of $X^{\ell-1}-1$ exists by Hensel's
lemma and is multiplicative by uniqueness. Thus $\omega^{\ell-1}=1$,
and its residual Galois character equals the reduction of the
$\ell$-adic cyclotomic character. On the common units take
$\varphi$ to be the primitive character inducing
$\omega^{2e}\theta^2$. If $C=\operatorname{cond}(\theta^2)$,
the product is already defined at level $\ell C$ by inflating
the two factors. Hence its conductor $f$ divides $\ell C$,
whether or not $\ell\mid C$. In particular $f>0$ is odd.
HXD1 shows that both factors are killed by $L$. The induced primitive
character is also killed by $L$, by surjectivity of reduction on unit
groups to each divisor of the modulus. Moreover $\varphi(-1)=1$
since both factors are squares at $-1$.

Let $\psi_2^\circ$ be the unique primitive character of conductor
$2^t$ inducing $\psi_2$. On units modulo $2^t$ take
$(\psi_2^\circ)^2\chi_4$, where $\chi_4$ is the nontrivial
character of $(\mathbb Z/4)^\times$ in $O^\times$. Its values on
units modulo $N$ agree with $\psi_2^2\chi_4$.
The host assumption $t\ge4$ gives
$(\mathbb Z/2^t)^\times=\langle-1\rangle\times\langle5\rangle$,
where $5$ has order $2^{t-2}$. The last assertion follows by
induction from $v_2(5^{2^j}-1)=j+2$, and the powers of $5$ are
precisely the classes $1\bmod4$. The kernel of reduction from
$2^t$ to $2^{t-1}$ is generated by $5^{2^{t-3}}$. Primitivity
at $2^t$ forces $\psi_2^\circ(5)$ to have exact order $2^{t-2}$.
Consequently $((\psi_2^\circ)^2\chi_4)(5)$ has exact order $2^{t-3}$,
as $\chi_4(5)=1$. The product descends to level $2^{t-1}$.
It cannot descend to $2^{t-2}$, where the order of the
class of $5$ is $2^{t-4}$, including $t=4$. This proves primitivity
of the resulting character $\xi$ at $2^n$. At $-1$ it takes value $-1$, giving
$\varphi(-1)\xi(-1)=-1$.

The input witness $q$ of the host supplies precisely $q>0$, $C\mid q$,
and $2(\ell q)4^{n_0(\ell,\psi)}\le 2^t$.
Since $f\mid\ell C\mid\ell q$, cancellation of $2$ gives
\[
       f\,4^{v_2(L)}\le(\ell q)4^{v_2(L)}\le2^{t-1}=2^n .
\]
Together with HXD1 this verifies every field of the Bernoulli
$\operatorname{XData}(\ell,d,\varphi,\xi)$, with the original
$\ell$, the prescribed $e$, and both characters valued in $O$.
% END HXD2

% BEGIN HXD3
\paragraph{HXD3. The identity on the absolute Galois group.}
Let $M=\ell N^2$. Since $C\mid N$ we have $f\mid\ell N$, and
$2^n\mid N$, so $f2^n\mid M$; also $4$, $\ell$, and $N$ divide $M$.
For $g\in G_{\mathbb Q}=\operatorname{Gal}(\Omega/\mathbb Q)$,
choose an integer $a$ representing its action on $\mu_M$. Then
$(a,M)=1$. Restricting the same action to $\mu_b$ for each $b\mid M$
gives $a\bmod b$. Thus evaluation of every character in the computation
below uses the same $a$. On these units the definition of the product
at the common level gives
\[
  (\operatorname{prodX}(\varphi,\xi))(a)=
        \big(\omega(a)^e\psi_2(a)\theta(a)\big)^2\chi_4(a).
\]
Write $\chi_\ell$ for the $O$-valued cyclotomic character and
$\psi_G$ for the Galois character induced by $\psi$.
By HXD1, $\bar\kappa(a)=1$. By HXD2,
$\overline{\omega(a)}=\bar\chi_\ell(g)$. The defining equations
for the supplied host frame give
$\Psi(\pi g)=\chi_\ell(g)^e\psi_G(g)$ and
$\eta(\pi g)=\chi_4(a)$. We obtain the equality of $k^\times$-valued
characters
\[
   \operatorname{redChar}(\operatorname{prodX}(\varphi,\xi))(g)
     =\bar\Psi(\pi g)^2\eta(\pi g).
\]
Here and below $\eta$ on the right is sent from $\{\pm1\}$ to
$k^\times$. In particular BR7 permits the inverse character
$\gamma(g)=\bar\Psi(\pi g)^{-2}\eta(\pi g)$, since
$\eta(\pi g)^{-1}=\eta(\pi g)$. Its coefficient field and reduction
map are $k$ and $O\twoheadrightarrow k$, and $f2^n>0$ as required.
% END HXD3

\section{Installing the residual vanishing theorem on the host frame}
\label{sec:hxp60}

Fix an arbitrary instance of the host problem. Thus $\Omega$ is the
specified algebraic closure of $\mathbb Q$, $\ell$ is a prime, $c$ is
the given coefficient datum, and $O=c.O$ has residue field $k$.
The supplied level $N\ne0$, exponent $e$, character
$\psi:(\mathbb Z/N)^\times\to O^\times$ and integer $o_1$ satisfy
the original $\mathrm{HostData}(\ell,e,\psi,o_1)$. Fix any $q$ satisfying
the host's $\mathrm{XData}(\ell,\psi,q)$. Write $G=\mathrm{Gal}(\Omega/\mathbb Q)$.
Take the supplied finite set $S$, its choices of valuation rings
$A_v$, the frame $F$, and all the identifications in
$\mathrm{IsHostFrame}$. In particular $S$ consists of the primes dividing
$\ell N$. The quotient $\Gamma=G_S$ kills the closed normal subgroup
generated by inertia above places outside $S$, and $F.\pi=\pi:G\to\Gamma$
is its quotient map. The distinguished $v_\ell$ lies over $\ell$.
Put $D_v=D_{A_v}$ and $I_v=I_{A_v}$. These are the actual subgroups
$F.D(v),F.I(v)$ of $G$, and $\mathrm{loc}_v$ is $\pi$ restricted to $D_v$.
The frame characters on $\Gamma$ are $\eta=F.\eta$ and $\Psi=F.\Psi$.
A bar on $\Psi$ denotes reduction $O^\times\to k^\times$.

% BEGIN HXP1
\paragraph{HXP1. The character and quotient for this frame.}
\label{claim:HXP1}
The scalar character needed for the weight $-2$ piece is
\[
   \lambda(s)=\bar\Psi(s)^{-2}\eta(s)\quad (s\in\Gamma).
\]
On $G$, $\gamma=\lambda\circ\pi$ is an inverse residual product
character covered by BR7 over $k$ with its given residue map.
The subgroup
\[
    K=\ker\big((\eta,\bar\Psi)\circ\pi\big)\trianglelefteq G
\]
lies in $\ker\gamma$ and has finite index prime to $\ell$.
Here and below an integer sign on the right of a field identity is
read in $k$ by its canonical map from $\mathbb Z$.

\begin{proof}
The ring $O$ is finite over $\mathbb Z_\ell$ and local with residue
characteristic $\ell$. Indeed the extension is integral, so
the contraction of its maximal ideal is the maximal ideal
$(\ell)$ of $\mathbb Z_\ell$. Consequently $k$ is a field quotient
of the finite-dimensional $\mathbb F_\ell$-algebra $O/\ell O$.
Write $|k|=\ell^r$ with $r\ge1$. The group $\eta$ takes values in
$\mathbb Z^\times=\{1,-1\}$, before reduction.

HXD3, applied with precisely the fixed $c,N,e,\psi,o_1,q$ and
$\Omega$, gives $d,f,n$, $\varphi$ of level $f$, and $\xi$ of level
$2^n$ over $O$. Their full Bernoulli character datum holds and
$f2^n\ne0$. For every $g\in G$ the residue of the Galois character of
$\operatorname{prodX}(\varphi,\xi)$ equals
\[
    \overline{\operatorname{Chr.Psi}(\ell,e,\operatorname{dirG}\psi)(g)}^{\,2}
       \,\eta_4(g)
      =\bar\Psi(\pi g)^2\eta(\pi g).
\]
The last equality uses both character identities of this host frame.
Since an integer sign equals its inverse, inverting this identity gives
$\gamma$ pointwise. We can therefore select the inverse alternative in
BR7 with $\mathrm{red}:O\to k$ the original residue map.
The signs and the residual units together have at most
$2(\ell^r-1)$ elements; more precisely the order of the image of
$(\eta,\bar\Psi)\circ\pi$ divides this number by Lagrange's theorem.
The host has $\ell\ne2$, hence the index of its kernel $K$ is prime
to $\ell$. Both character factors are one on $K$, which gives
$K\le\ker\gamma$.
\end{proof}
% END HXP1

% BEGIN HXP2
\paragraph{HXP2. The away-from-$\ell$ condition at every valuation.}
\label{claim:HXP2}
Let $z:\Gamma\to k$ be a locally constant $\lambda$-cocycle, and put
$u=z\circ\pi$. The frame's specified subgroup at $v$ is
\[
 N_v=\{s\in D_v : s\in I_v,\ \eta(\pi s)=1,\
                                     \bar\Psi(\pi s)=1\}.
\]
Suppose $z(\mathrm{loc}_v s)=0$ on $N_v$ at every $v\ne v_\ell$.
Then for every valuation ring $A$ of $\Omega$ with
$\operatorname{val}_A(\ell)\not<1$, the restriction of $u$ to
$I_A\cap K$ vanishes.

\begin{proof}
The quotient map is continuous, so $u$ is locally constant and
$u(gh)=u(g)+\gamma(g)u(h)$. In particular $u(1)=0$. For
$x\in\ker\gamma$, calculate
\[
   u(\sigma^{-1}x\sigma)
     =u(\sigma^{-1})+\gamma(\sigma)^{-1}
            \big(u(x)+u(\sigma)\big)
     =\gamma(\sigma)^{-1}u(x).                         \tag{HXP.1}
\]
This identity permits comparison with the one valuation above
each rational place chosen by the frame.

The valuation ring $A=\Omega$ has identity residue-field map, hence its
inertia is trivial. Every other $A$ lies above a rational prime $p$.
Indeed $A$ contains all algebraic integers, and its maximal ideal
contracts to a nonzero prime of $\mathbb Z$. If the contraction were
zero, every nonzero integer would be a unit in $A$; clearing an
integer denominator of an arbitrary algebraic number would then put
that number in $A$, forcing $A=\Omega$. The prime of contraction
determines the unique $v$ of $\mathbb Q$ below $A$. The hypothesis on
$\ell$ says $p\ne\ell$.

If $v\notin S$, the image of $I_A$ in $\Gamma$ is trivial by the
definition of $G_S$. Hence $u$ vanishes even on all of $I_A$.
If $v\in S$, the place differs from $v_\ell$. Conjugacy of valuations
above $v$ gives $\sigma\in G$ carrying $A_v$ to $A$. To recall
the infinite-extension point, on each finite normal subextension the
Galois group is transitive on the primes above $p$. The cosets of
lifts taking the first prime to the second are nonempty closed
subsets of the profinite $G$ and have the finite-intersection
property by passing to the compositum. An element in the intersection
carries the valuation rings to one another and conjugates their
inertia groups. For $x\in I_A\cap K$, it follows that
$s=\sigma^{-1}x\sigma\in I_{A_v}\cap K$ by normality of $K$.
Regarded as an element of $D_v$, $s$ belongs to the displayed $N_v$,
so $u(s)=z(\mathrm{loc}_v s)=0$. Equation (HXP.1) gives $u(x)=0$
since $K\le\ker\gamma$ and $\gamma(\sigma)$ is a unit. This also
handles every place in the host's special set away from $\ell$.
\end{proof}
% END HXP2

% BEGIN HXP3
\paragraph{HXP3. The exact tame-index condition and global scalar.}
\label{claim:HXP3}
Retain HXP2's local conditions on $z$ and suppose at the
distinguished place there exists $m_\ell\in k$ such that
\[
 z(\mathrm{loc}_{v_\ell}s)=\lambda(\mathrm{loc}_{v_\ell}s)m_\ell-m_\ell
                  \qquad(s\in D_{v_\ell}).
\]
Then some $m\in k$ satisfies
$z(s)=\lambda(s)m-m$ for every $s\in\Gamma$.

\begin{proof}
To apply BR7 on $G$, its original condition (BR.2) asks at every
valuation ring $A$ for a subgroup $H_A\le G$. With
$J_A=I_A\cap\ker\gamma$, the index $[J_A:J_A\cap H_A]$ must be
finite and prime to $\ell$, and $u$ must vanish on $J_A\cap H_A$.
If $\operatorname{val}_A(\ell)\not<1$, take $H_A=K$. The quotient
$J_A/(J_A\cap K)$ injects into $G/K$. Its order divides the
prime-to-$\ell$ order from HXP1. HXP2 supplies the required
vanishing, including for the trivial valuation.

Otherwise $A$ lies over $\ell$, as does $A_{v_\ell}$ by the
host's place equality. Choose $\sigma$ carrying $A_{v_\ell}$ to
$A$ as in HXP2. For $x\in J_A$, the element
$s=\sigma^{-1}x\sigma$ belongs to $I_{A_{v_\ell}}\subset D_{v_\ell}$
and $\gamma(s)=1$. The assumed local coboundary therefore gives
$u(s)=0$. Equation (HXP.1) gives $u(x)=0$. We may take $H_A=G$,
with relative index one. This proves (BR.2) for every $A$.
BR7 for the inverse character in HXP1 gives $m\in k$ with
$u(g)=\gamma(g)m-m$ for every $g\in G$. The quotient map is
surjective, and both sides are pullbacks from $\Gamma$. Evaluating at
any lift of $s\in\Gamma$ gives the asserted global formula there.
\end{proof}
% END HXP3

% BEGIN HXP4
\paragraph{HXP4. The weight $-2$ Selmer quotient on the CountV host.}
\label{claim:HXP4}
For every supplied datum and frame fixed above,
\[
    F.\mathrm{selmerCohomologyQuotient}(-2)=0 .
\]
Here zero means that every class is the zero class. This is
$F.\mathrm{Xpp}$, and the universal statement is precisely
$\mathrm{HostXpp}$.

\begin{proof}
Index the three rows and columns by $0,1,2$. The diagonal weights
are $1,0,-1$, so the matrix space $M$ of pure weight $-2$ consists
of the matrices $aE_{20}$, $a\in k$. In the base representation,
diagonal conjugation multiplies entry $(i,j)$ by
$\bar\Psi(s)^{\mathrm{wt}(i)-\mathrm{wt}(j)}$. When
$\eta(s)=-1$ the adjoint action first sends $X$ to $-JX^tJ$,
where $J$ reverses the three indices. The $(2,0)$ entry of this
new matrix is $-X_{20}$, because reversal and transposition fix
that position together. When $\eta(s)=1$ this step is the identity.
It follows that on $M$ the residual adjoint action is exactly
\[
    F.\bar g(-2)(s)(aE_{20})
              =\bar\Psi(s)^{-2}\eta(s)\,aE_{20}=\lambda(s)aE_{20}.
                                                               \tag{HXP.2}
\]

Let $b:\Gamma\to M$ represent any element of the actual Selmer
quotient. Its numerator consists of continuous cocycles whose
restriction to each $D_v$ belongs to the prescribed local submodule
$F.ZL(-2,v)$. Its denominator consists of those members of this
numerator which are global coboundaries. Since $M$ is discrete,
$z(s)=b(s)_{20}$ is locally constant. Taking that entry in the
cocycle identity, (HXP.2) gives
$z(st)=z(s)+\lambda(s)z(t)$. For $v\ne v_\ell$ the local submodule
requires $b(\mathrm{loc}_v g)=0$ for $g\in N_v$, so the scalar
equation of HXP2 holds. At $v_\ell$ the weight is strictly negative,
so the restriction of $b$ to the entire $D_{v_\ell}$ is a
coboundary. Taking the $(2,0)$ entry of the local witness gives
precisely the $m_\ell$ required in HXP3. The additional clause
in the local definition is indexed by weight zero and has false
antecedent here.

Apply HXP3 and set $B=mE_{20}\in M$. For every $s\in\Gamma$,
the matrix $F.\bar g(-2)(s)B-B$ and $b(s)$ have the same $(2,0)$
entry by (HXP.2); all their other entries are zero.
Hence $b$ belongs to the indicated denominator. This proves the
claimed vanishing for the original quotient over $k$.

In the CountV application a host $h$ supplies $o_1=h.o$, $q=h.q$
and $N=h.N$. The source extracts the two arithmetic predicates from
the character of this host, then chooses $S,A,F$ and the witness
of $\mathrm{IsHostFrame}$ once. Substituting exactly those objects
into the proved universal statement supplies the $F.\mathrm{Xpp}$
argument of $\mathrm{abstract\_count}$. The chosen marked element,
universal representation and length estimates are other arguments of
that application. This identifies the input at the original call site
without adding any hypotheses to $\mathrm{HostXpp}$.
\end{proof}
% END HXP4

\section{Arithmetic terms after the choice of a spectral setting}
Fix the original datum $D$, a setting $X$ and $Y:\mathrm{Full}(X)$.
For words $u,u'$ write $F,F'$ for their flags, and $U_F,U_{F'}$ for
the isotropic spans. All vector spaces and linear maps in this section
are over the original CM field $E$. An adelic term $s$ specifies two
row lattices and a double-unipotent class of rational isometries
$g_s:\Lambda'\longrightarrow\Lambda$. Choose the representative
specified in $s$, and write $G_s=g_sF'$ for its transported flag.
The map $s\mapsto t$ to arithmetic terms remembers the left lattice
and the double equivalence class of the transported frame. The finite
row index remembers the two lattices and a tuple of invertible matrices
$\iota_s$ at the level primes. These two maps forget different data.

% BEGIN ARF1
\paragraph{ARF1. Finiteness of each arithmetic-term fibre.}
Over any fixed arithmetic term $t$, there are finitely many adelic
terms. If the fibre is occupied, its size is at most the order of
the automorphism group of the projected positive definite lattice
of any one right row in the fibre.
\begin{proof}
Fix one lift with return $g$. For another lift with return $g'$,
the definition of the quotient of allowable frames gives two
integral unipotents $n\in N_{\Lambda,F}$, $m\in N_{\Lambda,gF'}$
such that the transported vectors of $g'$ are obtained by applying
$nm$ to those of $g$. Here equality of classes is double equivalence:
the relation is already transitive since conjugating an integral
flag unipotent by a lattice isometry transports its defining flag,
lattice and orthogonal-quotient conditions. Inverses give symmetry.
Equivalently
\[
       g'=n g\mu c,\qquad \mu=g^{-1}mg\in N_{\Lambda',F'},
       \qquad c=(nmg)^{-1}g'.
\]
The isometry $c$ fixes every right isotropic vector $e'_a$ and
the distinguished vector $z'$, and carries the right lattice of $g'$ to that of $g$. The two
right lattices are entries of the same list of representatives, and
$c$ is in the prescribed flag stabilizer. Uniqueness in the row
system therefore makes those two entries equal.

Choose such a $c$ for each lift. On $U_{F'}^\perp/U_{F'}$, it fixes
the negative vector $\bar z'$ and acts on its positive definite
orthogonal complement. It preserves the full projected lattice
in that complement. QPR5 proves finiteness of its isometry group
by taking the positive trace form at all infinite places.
If $c,d$ give the same action on the complement, then $c^{-1}d$
acts trivially on the entire quotient and fixes the flag vectors.
It preserves $\Lambda'$, so it is in $N_{\Lambda',F'}$ by definition.
Both $c$ and $d$ normalize this unipotent group. Hence the
corresponding returns are in the same double-unipotent class and
give the identical adelic term. The chosen assignment is injective.
An empty fibre is finite as well. This proves the field
\texttt{fibFin} for the original map $s\mapsto s.\mathrm{toTerm}$.
\end{proof}
% END ARF1

% BEGIN ARF2
\paragraph{ARF2. Recovering an action in characteristic zero from the row index.}
Let $s_i,t_i$ ($i=1,2$) have equal arithmetic terms within each
pair $(s_i,t_i)$. Suppose $s_1,s_2$ have equal finite row indices,
and likewise $t_1,t_2$. In the decompositions of ARF1 write
$g_{t_i}=n_i m_i g_i c_i$, with $g_i=g_{s_i}$. Then
$c_2c_1^{-1}$ acts identically on $U_{F'}^\perp/U_{F'}$.
\begin{proof}
All four left lattices agree, and all four right lattices agree,
by ARF1 and the row equalities. Use their base-column transports
$a,b$ for every term. Write $L(k)$ for the tuple of reduced
middle-block matrices, computed by pairing $k\,bw_j$ with $bwd_i$
and applying each of $Y.TR$'s residue maps. QPR8 proves that on
the integral group preserving the reference triangular $(e,z)$
congruences this operation is multiplicative and has invertible
values. Indeed the span of $e$ is invariant in the reduction and
these matrices give the action on the $bw$ quotient off that span.
This works for an arbitrary typed reduction by its stated kernel
and idempotent laws. In particular $\iota_s=L(a^{-1}g_{s,f}b)$,
including terms for which that return is outside the smaller level.
The base-column transport of a rational row unipotent has $L=1$:
the strengthened congruence in PGR3 and PGR9 places its displacement
of every $bw_j$ in the $e$-span modulo the modulus. More explicitly, set
$\mu_i=g_i^{-1}m_i g_i\in N_{\Lambda',F'}$. Then
\[
 a^{-1}g_{t_i,f}b=(a^{-1}n_{i,f}a)(a^{-1}g_{i,f}b)
                        (b^{-1}\mu_{i,f}b)(b^{-1}c_{i,f}b).
\]
The first and third factors have reduced block one. Thus
\[
               \iota_{t_i}=\iota_{s_i}\,L(b^{-1}c_{i,f} b).
\]
Here $c_i$ preserves the right lattice and fixes $e'_j,z'$.
After the base-column transport it is integral and fixes these
reference vectors modulo the modulus, so all displayed matrices
are in the stated domain of $L$.

Write $c_i=v_i h_i$ in the rational splitting $N_{u'}H_{u'}$
of PGR2. Its Levi component extends the induced isometry of
$U_{F'}^\perp/U_{F'}$ by fixing the hyperbolic pairs. If
$b=h_b x_{u'}$, conjugate first by $h_b$. The integral projection
of PGR4 splits the two finite factors into lattice-preserving
factors: $h_b$ normalizes the unipotent, and $x_{u'}$ preserves
the reference lattice. Thus the unipotent has trivial reduced
block by PGR3, and the displayed $L$ is the one for $h_i$.
Cancel the two equal $\iota_{s_i}$ from the left in the equal
$\iota_{t_i}$. Multiplicativity shows that the integral element
$b^{-1}(h_2h_1^{-1})_f b$ has middle block $1$. Its $(e,z)$
congruences and the typed kernel criterion give membership in
$K_2$: in each $bw$-column the difference from the identity lies
in the $e$-span modulo the modulus, exactly as in the last step
of QPR8. Now $h_2h_1^{-1}$ is a rational element of $H_{u'}$
in $h_b K_{H,u'} h_b^{-1}$. The original column freeness PGR5
forces it to be $1$. The unipotents $v_i$ act trivially on the
quotient, which proves the assertion.
\end{proof}
% END ARF2

% BEGIN ARF3
\paragraph{ARF3. The corank and scalar test descend together.}
At the four terms of ARF2, both outputs required by
\texttt{FinX.pairIndep} agree. They are the integer $s_0$ and
the truth value of the source's scalar test.
\begin{proof}
Write $U$ for the left isotropic span and $V_i=g_i U_{F'}$.
Denote the distinguished vectors by $z$ on the left and $z_i$
in $G_i=g_iF'$. The new spans and vectors are
$V_i'=n_i V_i$, $z_i'=n_i m_i z_i$.
In particular $z_i'-n_i z_i\in V_i'$, since
$z_i\in V_i^\perp$. Similarly $n_i z-z\in U$.
Set $\beta=g_2g_1^{-1}$ and
$\tau=g_1(c_2c_1^{-1})g_1^{-1}$. Then
\[
 \beta'=g_{t_2}g_{t_1}^{-1}
      =n_2m_2\beta\tau m_1^{-1}n_1^{-1}.
\]
The middle correction $\tau$ acts identically on
$V_1^\perp/V_1$ by ARF2. For $y\in V_1^\perp$, the four maps
$m_1^{-1},\tau,\beta,m_2$ induce in succession the identity,
identity, $\bar\beta$, identity on the indicated orthogonal
quotients. Thus
$\beta'n_1 y-n_2\beta y\in V_2'$. For $x_i\in U^\perp$
one also has $n_i x_i-x_i\in U$.

To apply these congruences to the actual test spaces, put
$B_i=U+Ez+V_i+Ez_i$. This is the source's $\mathrm{UUE}$:
intersecting the two rational isotropic spans with the full row
lattice before taking their $E$-span gives those same spans by
clearing denominators. The spans for the new pair are $B_i'=n_iB_i$,
as follows from the two vector congruences above. In particular
the orthogonals $B_i^\perp$ transform by $n_i$. The paired subspace
defining $s_0$ consists of the first coordinates of pairs satisfying
\[
 x_i\in B_i^\perp,\qquad x_1-x_2\in U,\qquad
                   \beta x_1-x_2\in V_2.
\]
The bijection $(x_1,x_2)\mapsto(n_1x_1,n_2x_2)$ transports all
three conditions in both directions by the congruences. The first
projections therefore have the same dimension. Subtracting it from
the prescribed $k$ gives the same $s_0$, including when that
natural-number subtraction is zero.

The scalar test uses $y_i\in B_i\cap(U+Ez)^\perp$ and
$y_i'\in B_i\cap(V_i+Ez_i)^\perp$. Their primed-setting spaces
are again the images by $n_i$. Its two compatibility hypotheses
are $y_1-y_2\in U$ and $y_2'-\beta y_1'\in V_2$.
Both are preserved and reflected by the same two congruences.
For every pair of scalars $a,b$, the relation tested on side $i$ is
\[
            y_i-y_i'\ \in\ az+bz_i+U+V_i .
\]
Applying $n_i$ gives exactly the relation for the new distinguished
vectors, with the very same $a,b$, because the differences in $z,z_i$
lie in the displayed spans. The equivalence between the two side
relations is therefore unchanged for every admissible quadruple
and all $a,b$. This is the full scalar clause, including the
degenerate pairs. Equality of rows ensures that the original
definition of $s_0$ uses its matching-lattices branch.
\end{proof}
% END ARF3

% BEGIN ARF4
\paragraph{ARF4. The real Gram parameter and the common signature scale.}
All arithmetic terms have $G\in F_{\rm real}$, the given totally
real subfield of $E$. For every signature $\epsilon$ on $n$
coordinates there is one positive scale for the archimedean
measure valid for every ordering of the positive and negative
coordinates. These are the fields \texttt{Greal} and \texttt{cdEx}.
\begin{proof}
If $D_j$ is the leading $j$-minor of the cross-pairing matrix of
the two framed configurations, their Gram parameter is
$D_{r+1}\overline{D_{r+1}}-TT'D_r\overline{D_r}$. Both $T,T'$
are the prescribed natural numbers. Conjugation fixes this difference
and its fixed subfield is exactly the original totally real field.

Count the positive and negative coordinates, $P$ and $Q$. Choose
a sign-compatible ordering and push the given Haar measure
$D.cd(P,Q)$ through the corresponding unitary matrix reindexing
onto $U(\epsilon)$. This continuous group isomorphism produces
a nonzero Haar measure, so Haar uniqueness gives a finite positive
multiple $c$ of the specified $\mu_G(\epsilon)$. A different
ordering has the same $P,Q$ and differs by a permutation
preserving signs. Reindexing by that permutation is an automorphism
of finite order. Its pushforward multiplies Haar measure by a
positive number whose finite power is one; the number is one.
Thus the same $c$ works for every ordering. On each measurable
set the pushforward is exactly the measure of its inverse image
under the matrix reindexing in \texttt{CdTieAt}. This includes
empty positive or negative blocks and the zero-dimensional group.
\end{proof}
% END ARF4

\section{The supplied Bessel wave and the choice of spectral data}

Fix an original datum $D$, a setting $X$ over $D$, and, where it is
needed, $Y:\mathrm{Full}(X)$. Two spaces occur. The space
$\mathcal H_X=\mathrm{HHX}(X)$ is the image of the finite and infinite
operators on the adelic automorphic quotient. The real Bessel quotient
for a word $u$ is $S_{u,\infty}\backslash G_\infty$, with sections
transforming on the left by $\psi_{u,\infty}\pi_u$. The Bessel wave
is a section on this second quotient. We will use a dimension bound
on each of the two spaces, for different purposes.

% BEGIN SFS1
\paragraph{SFS1. Finite dimension of the actual automorphic range.}
\label{claim:SFS1}
If the original infinite kernel $X.\mathrm{ph}$ is integrable,
$\dim_{\mathbb C}\mathcal H_X$ is finite. There is no gap condition
in this statement.

\begin{proof}
At each real place let $w_v$ be the given lowest unit vector,
$\sigma_v$ the given discrete series, $\mu_v$ the prescribed
group measure, and put
\[
 d_v=\left(\int|\langle w_v,\sigma_v(g)w_v\rangle|^2\,d\mu_v(g)
       \right)^{-1},\qquad
 p_v(g)=d_v\overline{\langle w_v,\sigma_v(g)w_v\rangle}.
\]
This is $X.\mathrm{dv}(v)$: the numerator in the source's
formal-degree formula is $\|w_v\|^4=1$. Discrete-series Schur
orthogonality gives a positive formal degree for $\mu_v$.
Specializing all four vectors in that identity to $w_v$ shows
that this degree equals $d_v$. Thus its four-vector identity holds
with precisely the scalar above. In particular $p_v$ is continuous
and $p_v(1)=d_v>0$.

By definition $X.\mathrm{ph}(g)=\prod_v p_v(g_v)$ and the
infinite measure is the product of the $\mu_v$. Fix one place $v$.
At every other place choose a relatively compact identity
neighborhood on which $|p_j|\ge d_j/2$.
Its Haar measure is finite and positive. Tonelli on the product
of those neighborhoods with the entire $v$-factor bounds a
positive constant times $\int |p_v|$ by $\int |X.\mathrm{ph}|$.
For a single place the constant is one. Hence every local kernel
has the $L^1$ property required in RFD8.

Apply RFD8 to $X.B$'s adelic domain, measure, right regular action,
finite test function, and finite convolution operator. Local
constancy and support in the lattice stabilizer are fields of its
admissible finite test. The local representations are irreducible
and unitary by the discrete-series field. The previous paragraph
supplies the kernel hypotheses. The weak infinite convolution
operator is exactly $X.R_{\mathrm{ph}}$ by its defining field.
The range in RFD8 is therefore $\mathcal H_X$ on the supplied quotient.
\end{proof}
% END SFS1

% BEGIN SFS2
\paragraph{SFS2. Positivity of both original word volumes.}
\label{claim:SFS2}
For every supplied $X,Y,u$ the level normalization holds and
\[
 0<Y.\mathrm{volN}(u)=\mu_f(K_2)_{\mathbb R}\,
                  \mu_{N,u}(Y.FN(u))_{\mathbb R},\qquad
 0<Y.\mathrm{volH}(u)=\mu_{H,u}(G_\infty)_{\mathbb R}.
\]
Here the real conversion of an infinite measure is zero by source
convention, so actual finiteness of each factor is part of the proof.

\begin{proof}
ASG14 applies to the reference frame of $X$: its residue
congruences place $K_g$ in $K_2$ with trivial character, and place
$K_2$ inside the lattice stabilizer. The former is open and the
latter compact. Since $K_2$ is a subgroup, it is compact open
and has positive finite mass for the given finite Haar measure.
These are the four original normalization clauses.

For $N=N_{u,\infty}$ use exactly the rational group
$L=\mathrm{NOinf}(X,u)$: at the finite places its elements
belong to $N_u(\mathbb A_f)\cap x_uK_2x_u^{-1}$. This is a
compact open subgroup of $N_u(\mathbb A_f)$. In particular $L$
contains all the integral exponentials needed for the compact
cover in FLD2. In detail, for a rational Malcev basis $Z_i$ as
in FLD2, clear the denominators of the exponential polynomials
and then multiply each $Z_i$ by a common multiple of the
congruence primes to the powers required by that finite open
subgroup. At the finitely many exceptional places the exponential
of a sufficiently divisible integral multiple approaches one
$p$-adically. At all other places integrality suffices. Both
$\exp(\pm Z_i)$ after these scalings belong to $L$. The decreasing
coordinate reduction of FLD2 gives a compact $C\subset N$
with $LC=N$. Also $L$ is countable, being an image of rational points.

The measure $\mu_{N,u}$ is supported on $N$ and restricts there
to a nonzero Haar measure finite on compact sets by the
sigma-finite invariance argument of FLD2. Put $A=Y.FN(u)\cap N$.
The supplied strict fundamental-domain field, restricted to $N$,
makes $A$ measurable with one representative per $L$-orbit.
If it were null then so would be $N=\bigcup_{l\in L}lA$.
To bound its measure above, enumerate $L$ and partition $A$ by
the first $l$ for which $la\in C$. The translated pieces are
disjoint by strictness, so $\mu_{N,u}(A)\le\mu_{N,u}(C)<\infty$.
This proves the assertion on the very domain stored in $Y$.
At infinity the residual $H_{u,\infty}$ is compact by PGR2 and
the definite orthogonal complement of PGR5. Its given measure
restricts to nonzero Haar measure by FLD2, hence its total mass
is finite and positive. Taking real values proves the formulas.
\end{proof}
% END SFS2

% BEGIN SFS3
\paragraph{SFS3. Multiplicity of the local projector on Bessel sections.}
\label{claim:SFS3}
Use any one of the actual real places and words under the original
$\mathrm{GapHyp}(D)$. Write $G_v$ for its real unitary group,
$S_v=N_v\rtimes H_v$ for its Bessel subgroup,
$\rho_v=\psi_v\pi_v$ for the given unitary action on $V_v$,
and $\mathcal I_v$ for the Hilbert space of sections with
$f(sx)=\rho_v(s)f(x)$ and finite squared norm on $S_v\backslash G_v$.
Convolution in the right translation action $R_v$ by $p_v$ of
SFS1 is an orthogonal projection of rank at most one on
$\mathcal I_v$.

\begin{proof}
BWP6, BWP20 and BWP23 give the integrability of $p_v$ for the
supplied Haar measure and lowest vector. RFD5 applies to every
strongly continuous unitary action. Here the quotient is unimodular,
and translations act unitarily; continuity follows by density of
continuous compactly supported sections. Write $P_v=R_v(p_v)$.
For $b\in\operatorname{ran}P_v$, RFD6 constructs the bounded
intertwiner $A_b$ from the Hilbert representation $\sigma_v$ into
$\mathcal I_v$ by $A_b(\sigma_v(g)w_v)=R_v(g)b$. The Schur
compression formula gives $\|A_b\|=\|b\|$. Thus evaluation at
$w_v$ in this construction distinguishes the intertwiners attached
to distinct vectors in the image.

Every bounded such intertwiner $A$ defines a continuous map on
the space of smooth vectors of the discrete series by
\[
                     B_A(a)=(Aa)(1)\in V_v.
\]
Here evaluation requires justification, because $Aa$ initially
denotes an $L^2$ class. A smooth vector for the action on
$\mathcal I_v$ has weak derivatives of every order in the
infinitesimal translation directions. In a relatively compact
local trivialization of the smooth homogeneous bundle, these
directions span the quotient tangent space by transitivity.
Solve for coordinate derivatives in terms of a spanning subfamily
on a smaller chart. Bundle transition and vertical action
coefficients are smooth, with bounded derivatives on compact
subcharts. Induction, with a compact cutoff, bounds the local
Sobolev norm of order $j$ by a constant times a finite sum of the
Hilbert norms of translation derivatives of orders at most $j$.
The invariant quotient measure has a smooth positive density in
this chart. Local Sobolev embedding for $j>\dim(S_v\backslash
G_v)/2$ gives a continuous evaluation; taking all $j$ gives a
unique smooth representative. Bounded intertwining commutes
with all these Hilbert derivatives, so $B_A$ is continuous for
the smooth-vector topology of $\sigma_v$. It obeys
$B_A(\sigma_v(s)a)=\rho_v(s)B_A(a)$. If $B_A=0$, evaluating
this construction on $\sigma_v(g)a$ gives $(Aa)(g)=0$ for
every $g$ and smooth $a$. Density then gives $A=0$.

For the resulting multiplicity bound we use the published
archimedean Bessel uniqueness theorem of D. Jiang, B. Sun and
C.-B. Zhu, \emph{Uniqueness of Bessel models: the archimedean
case}, Theorem A and its flag description in Section 2
(\href{https://arxiv.org/abs/0908.1728}{arXiv:0908.1728}).
In its notation the successive isotropic spaces are the spans
of the first $j$ vectors of the given $e$-frame; adjoining the
negative line generated by $z_u$ gives the last flag space.
The ambient signature is $(k+r,r+1)$, and the residual group is
the compact $U(k)$ on the positive complement. Each simple
additive quotient $L_j$ is a complex line. BWP20 computes on it
the additive character with nonzero coefficient $4\pi/g_D$,
including the last link to $z_u$. Its restriction to every
$L_j$ is therefore nontrivial, as genericity requires. It is unitary, and is trivial on the
compact part in the scalar-character formulation.
The representation on the large group is the smooth
globalization of the given irreducible discrete series, and the
second representation in the theorem is the contragredient
of the given irreducible compact type $V_v$. These are irreducible
Casselman--Wallach representations. The functionals to which the
theorem applies are
$(a,\ell)\mapsto\ell(B_A(a))$ on
$\sigma_v^\infty\widehat\otimes V_v^*$.
They transform by exactly $\psi_v$ since the two compact actions
cancel. The finite-dimensional dual separates the values of
$B_A$, and evaluation distinguished the intertwiners above.
Theorem A bounds this Hom space by dimension one. This proves
the assertion for the original coefficient projector. The
bound concerns its image on Bessel sections for this specified
representation and word.
\end{proof}
% END SFS3

% BEGIN SFS4
\paragraph{SFS4. Decay and compact reproduction for the supplied wave.}
\label{claim:SFS4}
The full original $\mathrm{WuDecayOfGap}$ holds at the published
multiplicity-one input of SFS3. For every $D$ satisfying the
original gap and \emph{every} supplied $X$ over it, the actual
functions $X.W_u$ satisfy weighted integrability and compact
reproduction for all words, and $X.\mathrm{ph}$ is integrable.

\begin{proof}
Fix the word. The real stabilizer, character and compact
representation on the full tensor space are respectively the
products of the local stabilizers, characters and compact
actions. For the representations this follows from the supplied
tensor inner-product, spanning and compact equivariance fields,
and the trivial action on the unipotent radical. PGR2 gives
the semidirect decomposition. BWP19 identifies the pushforward
of the given quotient measure to $S_\infty\backslash G_\infty$
with a finite positive multiple of the product of local quotient
Haar measures. Its proof uses a cutoff and uniqueness of the
quotient measure, so the identity extends from product test
functions to the whole Borel sigma-algebra. In particular the
Hilbert space of global Bessel sections is, up to an overall
norm factor, $\widehat\bigotimes_v\mathcal I_v$.
One may see surjectivity in measurable local trivializations:
product simple functions of finite support in measure and pure
fiber vectors span a dense subspace. This identification intertwines
the componentwise translations; the common scalar in the norm
does not change their integrated action.

By SFS3, $p_v\in L^1$ at every place. Finite-product Fubini gives
$X.\mathrm{ph}=\prod_v p_v\in L^1$ and the integrated
operator on sections equals $P=\widehat\bigotimes_v P_v$.
Indeed its weak matrix coefficients agree on pure tensors,
hence everywhere by the $L^1$ bound. Therefore
$\dim\operatorname{ran}P\le1$. The supplied function $W=X.W_u$
is already a normalized continuous equivariant section by
$X.Wu\_spec$. Its pointwise convolution by $X.\mathrm{ph}$
equals $W$. This equality also holds in the Hilbert space:
pair with an $L^2$ section $b$ and apply Fubini, using
$\int |X.\mathrm{ph}(g)|\,\|b\|_2\|R(g)W\|_2\,dg<\infty$.
Thus $PW=W$ and $\operatorname{ran}P$ is a line.

Use separately the packet $W^0$ of BWP20--BWP22, with nonvanishing
from BQS23 and the numerical implication BWP23. Their construction
uses the frames, actual compact representation, lowest vectors,
tensor map, and subgroup and quotient measures supplied in this
same $X$. In particular it constructs, for this comparison, a
normalized section for the identical $\rho_u$ and $P$.
Its defining convolution shows $PW^0=W^0$. Hence $W=cW^0$
in the quotient Hilbert space with $|c|=1$ by the common
normalization. Both sections are continuous, and their difference
has a continuous invariant norm. Haar quotient measure has full
support, so the equality holds for every point of the group.

The compactly supported bounded reproducing kernel constructed for
$W^0$ in BWP22 consequently reproduces $W$ pointwise, for the
same infinite Haar measure. Their absolute norms coincide, so
the weighted $L^1$ bound of BWP22 transfers with the literal
$\mathrm{nuInf}^{N_X}$ and $X.\nu_Q(u)$. The scalar may depend
on the word. This proves both clauses word by word together
with integrability of the common infinite kernel.
\end{proof}
% END SFS4

% BEGIN SFS5
\paragraph{SFS5. The full finiteness and compatibility record on $Y$.}
\label{claim:SFS5}
For any $D,X,Y$, the two hypotheses
$\mathrm{WuDecay}(X)$ and $X.\mathrm{ph}\in L^1$ imply
$\mathrm{FinX}(Y)$. In particular the record exists on every
supplied tied pair $\mathrm{ConcTie}(X,Y)$ without an arithmetic
gap hypothesis. Under $\mathrm{GapHyp}(D)$ it exists for every
$X,Y$. This gives the original $\mathrm{FinXC}$ for every
spectral tie over such a $D$.

\begin{proof}
For each word PGR6 identifies the original map
$\mathrm{sharp}$ with evaluation on the supplied finite column
representatives followed by Riesz duality. Its explicit inverse
assigns to a representative the corresponding functional and
extends by covariance; PGR5 makes this extension unique.
This supplies both injectivity and surjectivity, with their actual
complex-antilinear convention. SFS2 and FLD3 give folding on
every word. QPR3 uses $X.\mathrm{ph}\in L^1$ and those folds to
put the Poincar\'e functions in $\mathcal H_X$. For the particular
transversal $Y.Dt(u)$, the seed for a sum of two columns is the
sum of their seeds by PGR7 and the defining linear evaluation;
scalar multiplication has the same compatibility. Absolute
summability pointwise from the fold/range statement allows
addition and scalar multiplication through the period series.
Passing to their $L^2$ classes respects these operations. These
are the original $\mathrm{inHH}$ and $\mathrm{puLin}$.

SFS1 gives finite dimension of $\mathcal H_X$. ARF1 supplies
finiteness over each arithmetic term; ARF2--ARF3 give independence
of both the pair corank and its scalar predicate on equal-row
lifts. ARF4 puts every Gram parameter in the real fixed subfield
and supplies the one positive signature scaling for the Haar
comparison. SFS2 gives positive real volumes for both factors
at each word. These are all ten fields of $\mathrm{FinX}(Y)$.
The tied specialization takes the two analytic hypotheses from
the original $\mathrm{ConcTie}$. For the gap specialization
take them from SFS4 on the given $X$, and for
$\mathrm{FinXC}$ use $X=T.X$, $Y=T.Y$ for the given spectral tie.
\end{proof}
% END SFS5

% BEGIN SFS6
\paragraph{SFS6. The local comparisons and the concrete spectral package.}
\label{claim:SFS6}
Suppose $T:\mathrm{STies}(D)$ is supplied. If $T.X$ satisfies
the two analytic hypotheses of SFS5 and
\[
   1\le T_w\left(\frac{D.\mathrm{host}.\mathrm{ka}}
                                {\mathrm{gg}(D)}\right)^{2r}
            \quad\text{for every word }w,\qquad
            T_w=(D.\mathrm{mem}(w)).A.T,\quad r=D.R.rS ,
\tag{SFS.1}
\]
then the original $\mathrm{STiesOK}(T)$ holds and there is a
$\mathrm{ConcK}(D)$ with $X=T.X$, $Y=T.Y$. The original
$\mathrm{GapHyp}(D)$ suffices for the two analytic hypotheses.

\begin{proof}
Write $h_F$ for SFS5 and specialize the original field of $T$ to
this very witness to obtain $I=T.\mathrm{hI}(h_F)$, the
comparison $I:D.s\simeq s_Y$ with $s_Y=\mathrm{sOf}(h_F)$.
For every finite prime $p$ its local field $I.\mathrm{atP}(p)$
is an isomorphism from $(D.s).\mathrm{atP}(p)$ to
$s_Y.\mathrm{atP}(p)$. The latter is \emph{defined} to be
$\mathrm{locC}(D,p)$. This supplies the entire local comparison
$\mathrm{LocCP.LocTied}(D)$ using all primes, with no choices
of new local models.

The isometry $I.\mathrm{eCu}(w)$ identifies the given abstract
column module with $\mathrm{ColT}(h_F.\mathrm{shInj}(w))$.
The original column comparison $D.s\_Cu(w)$ identifies the rank
of the former with $D.\mathrm{duS}(w)$, so the real-valued
$\mathrm{duC}$ equality follows with these exact spaces.
The equality for $D_\ast$ is the mass-normalization field of
$T.\mathrm{autTie}$. SFS2 gives the normalized level and PGR5
gives neatness for all words. Use the original $T.\mathrm{hb}$,
$T.\psi$ and their supplied measure and character comparisons for
the archimedean and finite components; append the local isomorphisms
just obtained. The two analytic hypotheses and (SFS.1), which is
$\mathrm{TFrameAt}(D)$ verbatim, give the remaining clauses of
$\mathrm{ConcTie}(T.X,T.Y)$. Together these statements also
give each conjunct in $\mathrm{STiesOK}(T)$. The required
concrete package is $(T.X,T.Y,h_F,I,\mathrm{ConcTie})$.
Under the original gap SFS4 applies to the specific setting $T.X$.
In particular QPR11 now applies to every pair of its words, with
the given kernel normalization and the same comparison isometries.
\end{proof}
% END SFS6

% BEGIN SFS7
\paragraph{SFS7. Unconditional and hypothesis-dependent family thresholds.}
\label{claim:SFS7}
Let $\mathcal F$ be an original family with an $\mathrm{SSelF}$
witness $(a,T)$. If there is a number $b$ such that, for every
$x\ge b$ with $x>\mathcal F.x_1$ and $x\ge a$, every
$i\in\mathcal F.Ix(x)$ satisfies $\mathrm{GapHyp}(D_{x,i})$
and (SFS.1), then the family has an original $\mathrm{CSelF}$
witness. The same assertion holds with the two analytic
hypotheses for the selected $T_{x,i}.X$ in place of the gaps.
If this uniform-in-$i$ property holds on every family of a
class with an original $\mathrm{SSel}$, these choices give
$\mathrm{CSel}$ on that class.

There is also a conclusion when the numerical data are supplied
only under the original standing hypotheses. For any family class
with a supplied $\mathrm{SSel}$, a closed original
$\mathrm{PX.TFrame}$ event implies the closed $\mathrm{SelOK}$ event
at the parameter values
\[
 (C_g,c_b,C_u,c_T,C_T)
  =(\mathcal F.\mathrm{CYL}+1,\tfrac14,3,\tfrac12,\mathcal F.\mathrm{CTu}),
\]
for every late $x$ and every $i\in\mathcal F.Ix(x)$ in
$\mathrm{OwnEv}$, with the same $C_0$. It therefore implies the closed
$\mathrm{ConcEv}$ event. This latter implication
applies in particular to the selector induced by any supplied
$\mathrm{TSel}$ and its compatible $\mathrm{SelS}$ in
$\mathrm{link\_alt}$.

\begin{proof}
For the first assertion take $a'=\max(a,b)$. Given $x>\mathcal
F.x_1$, $x\ge a'$ and $i$, use the original $T_{x,i}$ at its
threshold $a$. SFS6 supplies its concrete package. This is the
required dependent function of $x$ and $i$, for all indices
simultaneously. If only existence of $b$ was supplied, choose it
once per family. If the original compatible member records come
from $\mathrm{TSel},\mathrm{SelS}$, project the latter to
$\mathrm{STies}$ at the former's threshold, exactly as in the
definition of the selector in $\mathrm{link\_alt}$. This preserves
both the chosen setting and the family point.

For the conditional event, strengthen the closed parameter
thresholds for $C',C_l$ to at least $29$ by taking their maxima
with the constant function $29$ at the middle threshold level.
A family large for the strengthened thresholds is still large
for the old ones, with the same outer accuracy bound. The exact
$\mathrm{OwnEv}$ uses the constants $sC'(\mathcal F)$ and
$sC_l(\mathcal F)$ fixed by this family. Under the event's
admissibility premise the standing inequalities give
$m\ge2$, $f\ge m^4\ge16$, $L>1$ and
$k\ge C'm^5 f L$. Hence $n=2r+k+1\ge1801$.
For adjacent weight indices, the same standing inequalities,
after cancelling their common shift, give
$b_{j+1}-b_j\ge C_l n^{19} f^5 L\ge16n^{19}$.
These are the original gap conditions. At the late threshold for
$\mathrm{PX.TFrame}$ the same admissibility premise, with
its conjuncts in the order used by $\mathrm{Cnt.Hyp0}$, gives
(SFS.1). Apply SFS6 to the selected $T_{x,i}$ for every $i$,
whenever its own threshold and the strict family threshold hold.
This is $\mathrm{SelOK}$. To obtain $\mathrm{ConcEv}$ enlarge
the event threshold to the maximum of its old value,
$a$, and $\mathcal F.x_1+1$. Above it the selected package
exists whenever the original admissibility premise holds.

The last implication is guarded by that premise; $\mathrm{CSelF}$
requires concrete packages at every late point even when the premise
fails. A common-family $\mathrm{SSel}$ and uniform arithmetic
frame inequalities across its points are retained at the scopes
stated here. The distinct-word, diagonal and counting estimates in
the full spectral comparison require further arguments.
\end{proof}
% END SFS7

\endgroup
